\documentclass[11pt]{article}
\usepackage{amsmath,amssymb}
\newcommand{\Gen}[1]{\langle #1\rangle}
\newcommand{\Mat}{\operatorname{M}}
\title{Certificates that an end generates the ambient group}
\author{}
\date{}
\begin{document}
\maketitle

\paragraph{Scope.}
The source census is Section 4 of Malcolm Hoong Wai Chen,
\emph{Abstract regular polytopes of $G_2(q)$ in even characteristic},
arXiv:2609.25004v1 (2026). Our finite domain uses its previously accepted
faithful natural model of the intended $G=G_2(11)$ and the accepted
catalogue, orbit and transporter identifications. The numbered positions
below belong to that verification domain. The argument supplies negative
intersection entries, not counterexamples to the source paper's census.

\paragraph{Borrowed classification.}
We use the maximal subgroup classification of P. B. Kleidman,
\emph{The maximal subgroups of the Chevalley groups $G_2(q)$ with $q$ odd,
the Ree groups ${}^2G_2(q)$, and their automorphism groups},
Journal of Algebra \textbf{117} (1988), 30--71,
doi:10.1016/0021-8693(88)90239-6.
Its specialization needed here is the following cover: a maximal subgroup
of $G$ is reducible on the natural seven dimensional module, or has one
of the abstract types
\[
 \operatorname{PGL}_2(11),\quad 2^3.\operatorname{L}_3(2),\quad
 G_2(2),\quad J_1,\quad \operatorname{PSL}_2(8),\quad
 \operatorname{PSL}_2(13).
\]
This conservative cover retains the last two types although their field
conditions exclude $q=11$. The reducible geometric families are the two
parabolics, $\operatorname{SL}_3(11):2$, $\operatorname{SU}_3(11):2$, and
$(\operatorname{SL}_2(11)\circ\operatorname{SL}_2(11)).2$.
Their invariant subspaces, respectively flags, lines, lines, and a
decomposition of dimensions $3+4$, are part of the borrowed natural-module
description. No classification or embedding theorem is claimed as new.
The accompanying source note records the accessible corroborating tables
and the precise limits of their inspection.

The invariant spaces can also be read directly in the octonion model.
Let $\mathbb O$ denote the normed octonion algebra, with unit $1$, and
put $V=1^\perp$. Its automorphisms fix $1$ and preserve the norm. The
parabolic embeddings stabilize an isotropic line or a distinguished
totally singular plane; the two $A_2$ normalizers stabilize nonsingular
lines in $V$. The orthogonal family preserves a nondegenerate quaternion
subalgebra $Q$. It therefore preserves
$(Q\cap V)\oplus Q^\perp$, whose dimensions are $3$ and $4$.
The identification of these geometric embeddings is standard; see
R. A. Wilson, \emph{The Finite Simple Groups}, Sections 4.3.5--4.3.6,
pp.~123--126, doi:10.1007/978-1-84800-988-2\_4.

\paragraph{Lemma.}
Let $H\leq G$. Suppose that $49$ actual words in $H$ are linearly
independent when regarded as matrices in $\Mat_7(\mathbb F_{11})$.
If $H$ contains an element of order $24$ or $44$, then $H=G$.

\paragraph{Proof.}
The $49$ matrices span all of $\Mat_7(\mathbb F_{11})$. If a maximal
subgroup containing $H$ preserved a nonzero proper subspace $W$, so would
every matrix in that span. The full matrix algebra does not preserve $W$:
given $0\ne v\in W$ and $u\notin W$, a linear map can send $v$ to $u$.
Thus only the listed abstract types remain.

In $\operatorname{PGL}_2(11)$, element orders divide $10$ or $12$, or are
$11$, according as the element is semisimple split, semisimple nonsplit,
or nontrivial unipotent. Neither $24$ nor $44$ occurs.
In $2^3.\operatorname{L}_3(2)$ an element of order $24$ would have image
of order $12$ or $24$ in $\operatorname{L}_3(2)$, since its intersection
with the elementary abelian kernel has order at most two. The latter
group has orders $1,2,3,4,7$, so this is impossible; the extension's
order $1344$ also excludes $44$.
The complete ATLAS class lists for $G_2(2)$ and $J_1$ contain neither
order. Finally $\operatorname{PSL}_2(8)$ has element orders $1,2,3,7,9$,
and $|\operatorname{PSL}_2(13)|=1092$ is divisible by neither $24$ nor
$44$. Each remaining maximal type is excluded. Hence $H=G$.
\hfill$\square$

The ATLAS data used in this proof are the author-maintained pages
\texttt{brauer.maths.qmul.ac.uk/Atlas/exc/G22/} and
\texttt{brauer.maths.qmul.ac.uk/Atlas/spor/J1/}. These element-order facts
are borrowed; the numerical certificates below do not recompute the ATLAS.

\paragraph{Application to a marking.}
For a marking $(r_0,r_1,r_2,r_3)$ in the fixed original order, let
$R=\Gen{r_1,r_2,r_3}$. A successful certificate for the lemma gives
$R=G$. The original-domain checks also give $r_0\ne1$. Consequently
\[
 \Gen{r_0}\cap\Gen{r_1,r_2,r_3}=\Gen{r_0}\ne1.
\]
The index sets $\{0\}$ and $\{1,2,3\}$ are disjoint. Their required
intersection subgroup in the intersection property is the identity.
Thus the marking fails the intersection property. There is no need to
enumerate a large end group or to infer word membership from membership
in a matrix algebra.

\paragraph{Certificate and verification.}
The supplied word tree starts at the identity; each other vertex is a
strictly earlier vertex multiplied by one of $r_1,r_2,r_3$. A separate
checker reconstructs all $49$ matrices and proves independence by
elimination over $\mathbb F_{11}$. A second supplied word $w$ is checked
by $w^{24}=1$, $w^{12}\ne1$, $w^8\ne1$, or by $w^{44}=1$,
$w^{22}\ne1$, $w^4\ne1$. Since the prime divisors are respectively
$2,3$ or $2,11$, these tests give the exact claimed order.

The accepted case list, certificate hashes, Source review, synthetic
controls, genuine calls, independent execution reviews and ledger adoption
are recorded separately. Until those records are complete, a discovered
word list alone earns no census credit. The shared argument is an
application of the borrowed classification; the work here supplies a
small certificate interface and the new finite case certificates. It
claims neither a new classification nor a solved open conjecture or a
Lean formalization.
\end{document}
